%%=============================================================================
%% Voorwoord
%%=============================================================================

\chapter*{\IfLanguageName{dutch}{Woord vooraf}{Preface}}%
\label{ch:voorwoord}

%% TODO:
%% Het voorwoord is het enige deel van de bachelorproef waar je vanuit je
%% eigen standpunt (``ik-vorm'') mag schrijven. Je kan hier bv. motiveren
%% waarom jij het onderwerp wil bespreken.
%% Vergeet ook niet te bedanken wie je geholpen/gesteund/... heeft

Ik zag AI vroeger als iets dat vooral in de wereld van sciencefiction films de verbeelding van de kijker deed dwalen in de richting van robots die zelf konden redeneren zoals mensen, maar geprogrammeerd waren om de instructies van hun creators te volgen. Ik dacht aan robots in films zoals Maria uit Metropolis(1927) en Johny 5 uit Short Circuit(1986), iets recenter, de griezelige zwerm van robots uit I, Robot(2004) en wie kan Haley Joel Osment's fantastische vertroebelen van de grens tussen mens en robot vergeten uit de film \textcite{A.I.2001} waar hij in de huid van een robot kruipt die zó goed als mens kan redeneren dat hij zelf ook echt gelooft een mens te zijn. Nu hebben we (nog) geen intelligente robots in ieder huis rondlopen... maar eigenlijk kwam die film toch verbazend dicht bij wat we ons nu inbeelden als AI met zijn alwetende computer-orakel "Dr. Know", volgens mij geïnspireerd door de wijze Wizard of Oz(uit het gelijknamig boek). Dit is ook volgens mij  hetgeen dat velen onder ons zich nu inbeelden als we vandaag denken aan AI. Een alwetend entiteit dat je via een scherm kan aanspreken en de antwoorden kent op elk mogelijke vraag. \\

Het concept van Artificiële Intelligentie heeft diepe roots en neemt verschillende vormen aan doorheen de tijd. In het algemeen, wanneer we denken aan het concept dat een niet levende machine (vb. computerprogramma) zelf complexe taken kan uitvoeren, zelf kan redeneren en zo steeds dichter bij komt in de buurt van hoe wij als mensen kunnen redeneren, wie zich in die richting verder blijft verplaatsen komt AGI(Artificial General Intelligence) tegen. Een concept dat voorlopig nog sciencefiction blijft. Volgens Alan Turing en gebruikmakend van zijn concept van de Turing Test, wanneer je als mens via text-interface interageert met een machine en het verschil niet kan waarnemen of je met een mens of met een machine praat, dan heeft deze de Turing Test overwonnen. De officiële benaming van AI zou later volgen, maar Turing's concept is vandaag de dag wel heel relevant waar we hetzelfde meemaken en de grens tussen echt en machine gefabriceerde tekst/beeld/geluid steeds moeilijker te determineren is... \\

Dit onderzoek was voor mij een groot ontdekkingsreis in de wereld van AI, ik heb ontzettend veel bijgeleerd en ook ontdekt dat deze kant van de software wereld ongelooflijk snel evolueert, op een tempo dat bijna niet bij te houden is, er worden regelmatig doorbraken gemaakt, nieuwe integraties en technieken ontwikkeld om AI in te zetten. Het is veelbesproken, (soms over-)hyped en één van de grootste topics voor discussie op dit moment in de wereld van software en digitale innovatie. Naast de hardware kantje van robotica zijn alle ontwikkelingen wat betreft AI ontzettend spannend voor mensen zoals ik die opgegroeid zijn als fan van sciencefiction verhalen. \\

\begin{comment}
    % j J
    %todo bedank wie me gesteund heeft
\end{comment}

